\subsection{Axes IR Declaration and Primitive Reference}
\label{app:axes-ir-declarations-primitives}
An Axes IR declaration binds ordered groups to operand positions and
specifies both output order and scalar semantics. The generalized dot below
pairs batch and contraction positions in list order; free axes preserve
operand order. Paired groups require matching extents, distinct in-range
selectors, and disjoint roles, including when groups are empty.

\begin{lstlisting}[language=TEPL,
  caption={Axes IR representation of generalized dot, with explicit attributes, axis bindings, and value semantics.},
  label={lst:dot-declaration}]
op dot(lhs: tensor, rhs: tensor) -> tensor {
  attrs {
    lb, lc, rb, rc: axis[];
  }
  axes {
    let lf = ordered_complement(all_axes(lhs), attrs.lb ++ attrs.lc);
    let rf = ordered_complement(all_axes(rhs), attrs.rb ++ attrs.rc);
    bind B { lhs: attrs.lb, rhs: attrs.rb };
    bind K { lhs: attrs.lc, rhs: attrs.rc };
    bind I { lhs: lf };
    bind J { rhs: rf };
    yield B ++ I ++ J;
  }
  semantics {
    yield Reduce(K, Product(lhs, rhs));
  }
}
\end{lstlisting}

\paragraph{Value and index expressions.}
\label{app:value-index-expressions}
The value primitives below compose operation semantics. Nonbroadcasting
addition binds corresponding operand axes and preserves their shared
order. Calculated reads use coordinate sequences, with one entry per
accessed axis. For example, a static slice reads at
\code{starts + strides * output_indices}. These coordinate calculations
are supported only for the modeled contracts and are outside the direct
whole-axis completeness theorem.
\begin{table}[!htbp]
\caption{Axes IR value and index primitives.}
\label{tab:axes-ir-value-index-primitives}
\small
\begin{tabularx}{\linewidth}{@{}lY@{}}\toprule
Primitive & Interpretation\\\midrule
\code{input}, \code{Read(input, indices)} & Operand read through axis bindings or explicit coordinates.\\
\code{Const(c)} & Signed integer scalar constant.\\
\code{Add(a,b)}, \code{Product(a,b)} & Scalar addition and multiplication.\\
\code{Reduce(G,e)} & Finite summation over group \code{G}; other reductions are not implied.\\
\code{Coords(G)} & Current coordinates in group order.\\
\code{IndexAdd(a,b)}, \code{Concat(a,b)} & Elementwise index addition and sequence concatenation.\\
\code{IndexScale(s,i)}, \code{IndexOffset(o,i)} & Entrywise scaling and offset of index sequence \code{i}.\\\bottomrule
\end{tabularx}
\end{table}

\subsection{Symbolic Mapping: Algorithm and Reassociation Example}
\label{app:symbolic-mapping-reassociation}
The mapping state is $\Sigma=(V,\mathcal M,\Pi)$, where
$V[t]=(\operatorname{Ord}_t,e_t)$ stores an ordered interface and a value expression.
A mapping $(\operatorname{Ord}_t,\vec G,q)\in\mathcal M$ satisfies
\begin{equation}
 \gath(\operatorname{Ord}_t,q)=\vec G.
 \label{extended:eq:mapping-invariant}
\end{equation}
Each partition in $\Pi$ lists pairs $(\vec G_i,q_i)$ that cover an
interface exactly once: the concatenated position lists are a permutation
of its positions. Different partitions may group the same interface
differently without changing the tensor. Algorithm~\ref{alg:mapping}
visits producers before consumers and maintains these two invariants.

\begin{paperalgorithm}[H]
\caption{Build interfaces, value expressions, and axis mappings.}
\label{alg:mapping}
\begin{lstlisting}[language=Python]
def build_symbolic_axis_mapping(graph):
    V = initialize_inputs(graph)
    mappings, partitions = set(), {}
    for op in producer_first(graph):
        inputs = [V[t] for t in op.inputs]
        groups, bindings = resolve_bindings(op, inputs)
        aligned = []
        for (axes, value), binding in zip(inputs, bindings):
            aligned_axes = build_interface(axes, binding)
            register_mappings(aligned_axes, binding, mappings)
            refine_and_register(aligned_axes, binding, mappings, partitions)
            value = transport(value, axes, aligned_axes, mappings, partitions)
            aligned.append((aligned_axes, value))
        out_axes = evaluate(op.output_axes, groups)
        binding = output_positions(op, groups)
        register_mappings(out_axes, binding, mappings)
        refine_and_register(out_axes, binding, mappings, partitions)
        V[op] = (out_axes, compose(op.value, aligned, groups))
    return V, mappings, partitions
\end{lstlisting}
\end{paperalgorithm}

\paragraph{Reconstruction and composition.}
For a complete binding, concatenate its groups and their position lists.
If $q=q_1\Vert\cdots\Vert q_m$ is a permutation, then
\begin{equation}
 Q=\gath(\vec G_1\Vert\cdots\Vert\vec G_m,\operatorname{inverse}(q))
 \label{extended:eq:rebuilt-interfaces}
\end{equation}
restores tensor dimension order, including interleaved groups.
\code{register_mappings} records the resulting selections;
\code{refine_and_register} compares partitions through paired overlap
matching. \code{transport} expresses reads in the aligned interface, and
\code{compose} applies the declared value semantics. Missing mappings are
queried by composing existing selectors with
Equation~\eqref{eq:gather-composition}. These operations preserve the
mapping invariant; acceptance still requires checking the completed rule.

\paragraph{Partitioning a contraction.}
\label{extended:sec:partitioning}
Use the guarded reassociation and notation of
Appendix~\ref{app:batched-reassociation-program}. Its source intermediate has
interface $\vec B\Vert\vec I\Vert\vec J$, and the guard preserves
$\vec B$ as the outer batch prefix. The outer contraction group
$\vec D=\gath(\operatorname{Ord}_T,d_\ell)$ therefore selects from $\vec I,\vec J$.
For either producer group $\vec G$ at positions $q_G$, compute
\begin{equation}
\begin{aligned}
 (u_G,v_G)&=\mpo(d_\ell,q_G),&
 z_G&=\gath(d_r,u_G),\\
 \vec G_{\mathrm{sel}}&=\gath(\vec G,v_G),&
 \vec G_{\mathrm{rem}}&=\gath(\vec G,\ocomp(\operatorname{range}(|\vec G|),v_G)).
\end{aligned}
 \label{extended:eq:z-partners}
\end{equation}
Here $u_G$ indexes contraction-list slots, while $v_G$ indexes axes
within the producer group. Only $u_G$ selects the corresponding partners
in $d_r$. For example, let $\operatorname{Ord}_T=[i_0,i_1,j_0,j_1]$,
$d_\ell=[3,0]$, and $d_r=[2,0]$. The partitions give
$(u_I,v_I,z_I)=([1],[0],[0])$ and
$(u_J,v_J,z_J)=([0],[1],[2])$: $i_0$ pairs with position zero in $Z$,
and $j_1$ with position two. Selected groups follow contraction order;
remaining groups preserve producer order. Their concatenation need not
reproduce the original group, so mappings retain both selectors.

\paragraph{From mappings to attribute programs.}
\label{app:mappings-to-attribute-programs}
For admitted plans, the symbolic adapter consumes the shared metadata IR
of Theorem~\ref{thm:symbolic-compilation}. Symbolic lowering first matches source reads to
the RHS arithmetic skeleton; unresolved reduction information is not an
assertion that no reduction occurs. Read correspondence supplies logical
axes for construction, and interprets the already emitted retained-set, ordering, and position
expressions through these mappings; it does not reconstruct their recurrences.
Restricted addition treats matched components independently.

In reassociation, $U=\operatorname{dot}(Y,Z)$ contracts $\Jsel$ but keeps
$\vec K$ and $\Isel$ for the final dot. Its actual interface, rather than
an assumed contiguous grouping, locates both retained groups.
Substituting that interface into the final position queries yields the
source-only expressions in Table~\ref{tab:reassociation-program}.
Supplied extent and coordinate expressions are preserved. Unresolved
queries or ambiguous read assignments leave symbolic lowering unsuccessful;
the shared residual program remains available for native interpretation.
The separate validation calculus checks definedness, legality, ordered
outputs, and complete values; its normalization laws are specified and
proved in Appendix~\ref{app:validation-proofs}. The full reassociation
derivation is in Appendix~\ref{app:reassociation-proof}.

Detailed alternative examples, expanded mapping tables, and tutorial
figures are retained in \nolinkurl{evaluation/supplementary-materials/}.
The shared construction correspondence and guard arguments below do not
require this optional symbolic lowering to simplify every query. Broader
nonadmitted candidates use the general constructor and independent validation.
